% =====================================================================
%  Gemeinsame Farben und Stile der Zeichnungen (TikZ, selbst erstellt)
%  In der Präambel einbinden: % =====================================================================
%  Gemeinsame Farben und Stile der Zeichnungen (TikZ, selbst erstellt)
%  In der Präambel einbinden: % =====================================================================
%  Gemeinsame Farben und Stile der Zeichnungen (TikZ, selbst erstellt)
%  In der Präambel einbinden: % =====================================================================
%  Gemeinsame Farben und Stile der Zeichnungen (TikZ, selbst erstellt)
%  In der Präambel einbinden: \input{zeichnungen/stil}
%  Angelehnt an die Zeichnungen der anderen FabLab-Einweisungen
%  (z.B. fau-fablab/shaper-origin-einweisung). Lizenz: CC BY-SA 3.0
% =====================================================================
\usetikzlibrary{calc,arrows.meta,decorations.pathreplacing,patterns}

\definecolor{okgruen}{RGB}{46,160,67}
\definecolor{nokrot}{RGB}{207,34,46}
\definecolor{mittel}{RGB}{230,150,0}
\definecolor{feld}{RGB}{244,246,248}
\definecolor{holz}{RGB}{226,196,150}
\definecolor{holzrand}{RGB}{150,110,60}
\definecolor{gehaeusegruen}{RGB}{88,110,80}   % Gehäuse der Säge
\definecolor{gelbteil}{RGB}{226,200,40}       % gelbe Bedienteile
\definecolor{schutzorange}{RGB}{240,120,40}   % Sägeblattschutz
\definecolor{alu}{RGB}{205,210,215}           % Sägetisch
\definecolor{gefahr}{RGB}{207,34,46}

\tikzset{
  % Bauteile der Maschine
  gehaeuse/.style={fill=gehaeusegruen!70, draw=black!80, line width=0.7pt, line join=round},
  tisch/.style={fill=alu, draw=black!75, line width=0.7pt, line join=round},
  metall/.style={fill=black!30, draw=black!75, line width=0.6pt, line join=round},
  gelb/.style={fill=gelbteil, draw=black!75, line width=0.6pt, line join=round},
  blatt/.style={fill=black!55, draw=black!85, line width=0.5pt, line join=round},
  keil/.style={fill=black!20, draw=black!85, line width=0.6pt, line join=round},
  werkstueck/.style={fill=holz, draw=holzrand, line width=0.6pt, line join=round},
  hand/.style={fill=orange!20, draw=orange!60!black, line width=0.6pt, line join=round, rounded corners=1.2mm},
  gefahrenzone/.style={pattern=north east lines, pattern color=gefahr!70, draw=gefahr, dashed, line width=0.5pt},
  % Pfeile, Maße, Beschriftung
  vorschub/.style={-{Stealth[length=2.6mm, width=2.2mm]}, line width=1.3pt},
  drehung/.style={-{Stealth[length=1.8mm, width=1.6mm]}, line width=0.9pt},
  kraft/.style={-{Stealth[length=2.2mm, width=2mm]}, line width=1.1pt},
  mass/.style={{Stealth[length=1.4mm]}-{Stealth[length=1.4mm]}, line width=0.4pt, black!80},
  hilfslinie/.style={line width=0.3pt, black!60},
  marke/.style={circle, draw=black, fill=white, line width=0.8pt,
    minimum size=4.6mm, inner sep=0pt, font=\scriptsize\bfseries},
  zeiger/.style={line width=0.6pt},
  % Bewertung: \pic at (x,y) {ok}; bzw. {nok}
  ok/.pic={\fill[okgruen] (0,0) circle (3);
    \draw[white, line width=1.3pt, line cap=round, line join=round]
      (-1.4,0.1) -- (-0.4,-1.1) -- (1.5,1.2);},
  nok/.pic={\fill[nokrot] (0,0) circle (3);
    \draw[white, line width=1.3pt, line cap=round]
      (-1.1,-1.1) -- (1.1,1.1) (-1.1,1.1) -- (1.1,-1.1);},
}

% Feld mit hellem Hintergrund und Überschrift: \feld{x}{y}{Breite}{Höhe}{Titel}
\newcommand{\feld}[5]{%
  \fill[feld, rounded corners=2mm] (#1,#2) rectangle ++(#3,#4);
  \node[anchor=north west, font=\small\bfseries] at (#1+3,#2+#4-2) {#5};}

% Nummer an einer Zeichnung mit Zeigerlinie: \nr{Nummer}{Position}{Ziel}
\newcommand{\nr}[3]{%
  \draw[zeiger] (#2) -- (#3);
  \fill (#3) circle[radius=0.6];
  \node[marke] at (#2) {#1};}

% Kreissägeblatt in der Seitenansicht: \saegeblatt{x}{y}{Radius}{Zähne}
% Die Zähne zeigen im Uhrzeigersinn. In allen Seitenansichten steht die
% Bedienperson rechts und schiebt das Werkstück nach links: Die Zähne
% laufen oben auf sie zu und schneiden vorn von oben nach unten.
\newcommand{\saegeblatt}[4]{%
  \begin{scope}[shift={(#1,#2)}]
    \pgfmathsetmacro{\sbs}{360/#4}%
    \pgfmathsetmacro{\sbz}{min(0.18*#3, 2*pi*#3/#4*0.55)}%
    \fill[black!60, draw=black!85, line width=0.4pt] (0:#3)
      \foreach \i in {1,...,#4}{
        -- ({\i*\sbs-\sbs}:#3) -- ({\i*\sbs-0.2*\sbs}:#3-\sbz) -- ({\i*\sbs}:#3)}
      -- cycle;
    \fill[black!25] (0,0) circle (#3*0.5);
    \fill[white] (0,0) circle (#3*0.12);
  \end{scope}}

%  Angelehnt an die Zeichnungen der anderen FabLab-Einweisungen
%  (z.B. fau-fablab/shaper-origin-einweisung). Lizenz: CC BY-SA 3.0
% =====================================================================
\usetikzlibrary{calc,arrows.meta,decorations.pathreplacing,patterns}

\definecolor{okgruen}{RGB}{46,160,67}
\definecolor{nokrot}{RGB}{207,34,46}
\definecolor{mittel}{RGB}{230,150,0}
\definecolor{feld}{RGB}{244,246,248}
\definecolor{holz}{RGB}{226,196,150}
\definecolor{holzrand}{RGB}{150,110,60}
\definecolor{gehaeusegruen}{RGB}{88,110,80}   % Gehäuse der Säge
\definecolor{gelbteil}{RGB}{226,200,40}       % gelbe Bedienteile
\definecolor{schutzorange}{RGB}{240,120,40}   % Sägeblattschutz
\definecolor{alu}{RGB}{205,210,215}           % Sägetisch
\definecolor{gefahr}{RGB}{207,34,46}

\tikzset{
  % Bauteile der Maschine
  gehaeuse/.style={fill=gehaeusegruen!70, draw=black!80, line width=0.7pt, line join=round},
  tisch/.style={fill=alu, draw=black!75, line width=0.7pt, line join=round},
  metall/.style={fill=black!30, draw=black!75, line width=0.6pt, line join=round},
  gelb/.style={fill=gelbteil, draw=black!75, line width=0.6pt, line join=round},
  blatt/.style={fill=black!55, draw=black!85, line width=0.5pt, line join=round},
  keil/.style={fill=black!20, draw=black!85, line width=0.6pt, line join=round},
  werkstueck/.style={fill=holz, draw=holzrand, line width=0.6pt, line join=round},
  hand/.style={fill=orange!20, draw=orange!60!black, line width=0.6pt, line join=round, rounded corners=1.2mm},
  gefahrenzone/.style={pattern=north east lines, pattern color=gefahr!70, draw=gefahr, dashed, line width=0.5pt},
  % Pfeile, Maße, Beschriftung
  vorschub/.style={-{Stealth[length=2.6mm, width=2.2mm]}, line width=1.3pt},
  drehung/.style={-{Stealth[length=1.8mm, width=1.6mm]}, line width=0.9pt},
  kraft/.style={-{Stealth[length=2.2mm, width=2mm]}, line width=1.1pt},
  mass/.style={{Stealth[length=1.4mm]}-{Stealth[length=1.4mm]}, line width=0.4pt, black!80},
  hilfslinie/.style={line width=0.3pt, black!60},
  marke/.style={circle, draw=black, fill=white, line width=0.8pt,
    minimum size=4.6mm, inner sep=0pt, font=\scriptsize\bfseries},
  zeiger/.style={line width=0.6pt},
  % Bewertung: \pic at (x,y) {ok}; bzw. {nok}
  ok/.pic={\fill[okgruen] (0,0) circle (3);
    \draw[white, line width=1.3pt, line cap=round, line join=round]
      (-1.4,0.1) -- (-0.4,-1.1) -- (1.5,1.2);},
  nok/.pic={\fill[nokrot] (0,0) circle (3);
    \draw[white, line width=1.3pt, line cap=round]
      (-1.1,-1.1) -- (1.1,1.1) (-1.1,1.1) -- (1.1,-1.1);},
}

% Feld mit hellem Hintergrund und Überschrift: \feld{x}{y}{Breite}{Höhe}{Titel}
\newcommand{\feld}[5]{%
  \fill[feld, rounded corners=2mm] (#1,#2) rectangle ++(#3,#4);
  \node[anchor=north west, font=\small\bfseries] at (#1+3,#2+#4-2) {#5};}

% Nummer an einer Zeichnung mit Zeigerlinie: \nr{Nummer}{Position}{Ziel}
\newcommand{\nr}[3]{%
  \draw[zeiger] (#2) -- (#3);
  \fill (#3) circle[radius=0.6];
  \node[marke] at (#2) {#1};}

% Kreissägeblatt in der Seitenansicht: \saegeblatt{x}{y}{Radius}{Zähne}
% Die Zähne zeigen im Uhrzeigersinn. In allen Seitenansichten steht die
% Bedienperson rechts und schiebt das Werkstück nach links: Die Zähne
% laufen oben auf sie zu und schneiden vorn von oben nach unten.
\newcommand{\saegeblatt}[4]{%
  \begin{scope}[shift={(#1,#2)}]
    \pgfmathsetmacro{\sbs}{360/#4}%
    \pgfmathsetmacro{\sbz}{min(0.18*#3, 2*pi*#3/#4*0.55)}%
    \fill[black!60, draw=black!85, line width=0.4pt] (0:#3)
      \foreach \i in {1,...,#4}{
        -- ({\i*\sbs-\sbs}:#3) -- ({\i*\sbs-0.2*\sbs}:#3-\sbz) -- ({\i*\sbs}:#3)}
      -- cycle;
    \fill[black!25] (0,0) circle (#3*0.5);
    \fill[white] (0,0) circle (#3*0.12);
  \end{scope}}

%  Angelehnt an die Zeichnungen der anderen FabLab-Einweisungen
%  (z.B. fau-fablab/shaper-origin-einweisung). Lizenz: CC BY-SA 3.0
% =====================================================================
\usetikzlibrary{calc,arrows.meta,decorations.pathreplacing,patterns}

\definecolor{okgruen}{RGB}{46,160,67}
\definecolor{nokrot}{RGB}{207,34,46}
\definecolor{mittel}{RGB}{230,150,0}
\definecolor{feld}{RGB}{244,246,248}
\definecolor{holz}{RGB}{226,196,150}
\definecolor{holzrand}{RGB}{150,110,60}
\definecolor{gehaeusegruen}{RGB}{88,110,80}   % Gehäuse der Säge
\definecolor{gelbteil}{RGB}{226,200,40}       % gelbe Bedienteile
\definecolor{schutzorange}{RGB}{240,120,40}   % Sägeblattschutz
\definecolor{alu}{RGB}{205,210,215}           % Sägetisch
\definecolor{gefahr}{RGB}{207,34,46}

\tikzset{
  % Bauteile der Maschine
  gehaeuse/.style={fill=gehaeusegruen!70, draw=black!80, line width=0.7pt, line join=round},
  tisch/.style={fill=alu, draw=black!75, line width=0.7pt, line join=round},
  metall/.style={fill=black!30, draw=black!75, line width=0.6pt, line join=round},
  gelb/.style={fill=gelbteil, draw=black!75, line width=0.6pt, line join=round},
  blatt/.style={fill=black!55, draw=black!85, line width=0.5pt, line join=round},
  keil/.style={fill=black!20, draw=black!85, line width=0.6pt, line join=round},
  werkstueck/.style={fill=holz, draw=holzrand, line width=0.6pt, line join=round},
  hand/.style={fill=orange!20, draw=orange!60!black, line width=0.6pt, line join=round, rounded corners=1.2mm},
  gefahrenzone/.style={pattern=north east lines, pattern color=gefahr!70, draw=gefahr, dashed, line width=0.5pt},
  % Pfeile, Maße, Beschriftung
  vorschub/.style={-{Stealth[length=2.6mm, width=2.2mm]}, line width=1.3pt},
  drehung/.style={-{Stealth[length=1.8mm, width=1.6mm]}, line width=0.9pt},
  kraft/.style={-{Stealth[length=2.2mm, width=2mm]}, line width=1.1pt},
  mass/.style={{Stealth[length=1.4mm]}-{Stealth[length=1.4mm]}, line width=0.4pt, black!80},
  hilfslinie/.style={line width=0.3pt, black!60},
  marke/.style={circle, draw=black, fill=white, line width=0.8pt,
    minimum size=4.6mm, inner sep=0pt, font=\scriptsize\bfseries},
  zeiger/.style={line width=0.6pt},
  % Bewertung: \pic at (x,y) {ok}; bzw. {nok}
  ok/.pic={\fill[okgruen] (0,0) circle (3);
    \draw[white, line width=1.3pt, line cap=round, line join=round]
      (-1.4,0.1) -- (-0.4,-1.1) -- (1.5,1.2);},
  nok/.pic={\fill[nokrot] (0,0) circle (3);
    \draw[white, line width=1.3pt, line cap=round]
      (-1.1,-1.1) -- (1.1,1.1) (-1.1,1.1) -- (1.1,-1.1);},
}

% Feld mit hellem Hintergrund und Überschrift: \feld{x}{y}{Breite}{Höhe}{Titel}
\newcommand{\feld}[5]{%
  \fill[feld, rounded corners=2mm] (#1,#2) rectangle ++(#3,#4);
  \node[anchor=north west, font=\small\bfseries] at (#1+3,#2+#4-2) {#5};}

% Nummer an einer Zeichnung mit Zeigerlinie: \nr{Nummer}{Position}{Ziel}
\newcommand{\nr}[3]{%
  \draw[zeiger] (#2) -- (#3);
  \fill (#3) circle[radius=0.6];
  \node[marke] at (#2) {#1};}

% Kreissägeblatt in der Seitenansicht: \saegeblatt{x}{y}{Radius}{Zähne}
% Die Zähne zeigen im Uhrzeigersinn. In allen Seitenansichten steht die
% Bedienperson rechts und schiebt das Werkstück nach links: Die Zähne
% laufen oben auf sie zu und schneiden vorn von oben nach unten.
\newcommand{\saegeblatt}[4]{%
  \begin{scope}[shift={(#1,#2)}]
    \pgfmathsetmacro{\sbs}{360/#4}%
    \pgfmathsetmacro{\sbz}{min(0.18*#3, 2*pi*#3/#4*0.55)}%
    \fill[black!60, draw=black!85, line width=0.4pt] (0:#3)
      \foreach \i in {1,...,#4}{
        -- ({\i*\sbs-\sbs}:#3) -- ({\i*\sbs-0.2*\sbs}:#3-\sbz) -- ({\i*\sbs}:#3)}
      -- cycle;
    \fill[black!25] (0,0) circle (#3*0.5);
    \fill[white] (0,0) circle (#3*0.12);
  \end{scope}}

%  Angelehnt an die Zeichnungen der anderen FabLab-Einweisungen
%  (z.B. fau-fablab/shaper-origin-einweisung). Lizenz: CC BY-SA 3.0
% =====================================================================
\usetikzlibrary{calc,arrows.meta,decorations.pathreplacing,patterns}

\definecolor{okgruen}{RGB}{46,160,67}
\definecolor{nokrot}{RGB}{207,34,46}
\definecolor{mittel}{RGB}{230,150,0}
\definecolor{feld}{RGB}{244,246,248}
\definecolor{holz}{RGB}{226,196,150}
\definecolor{holzrand}{RGB}{150,110,60}
\definecolor{gehaeusegruen}{RGB}{88,110,80}   % Gehäuse der Säge
\definecolor{gelbteil}{RGB}{226,200,40}       % gelbe Bedienteile
\definecolor{schutzorange}{RGB}{240,120,40}   % Sägeblattschutz
\definecolor{alu}{RGB}{205,210,215}           % Sägetisch
\definecolor{gefahr}{RGB}{207,34,46}

\tikzset{
  % Bauteile der Maschine
  gehaeuse/.style={fill=gehaeusegruen!70, draw=black!80, line width=0.7pt, line join=round},
  tisch/.style={fill=alu, draw=black!75, line width=0.7pt, line join=round},
  metall/.style={fill=black!30, draw=black!75, line width=0.6pt, line join=round},
  gelb/.style={fill=gelbteil, draw=black!75, line width=0.6pt, line join=round},
  blatt/.style={fill=black!55, draw=black!85, line width=0.5pt, line join=round},
  keil/.style={fill=black!20, draw=black!85, line width=0.6pt, line join=round},
  werkstueck/.style={fill=holz, draw=holzrand, line width=0.6pt, line join=round},
  hand/.style={fill=orange!20, draw=orange!60!black, line width=0.6pt, line join=round, rounded corners=1.2mm},
  gefahrenzone/.style={pattern=north east lines, pattern color=gefahr!70, draw=gefahr, dashed, line width=0.5pt},
  % Pfeile, Maße, Beschriftung
  vorschub/.style={-{Stealth[length=2.6mm, width=2.2mm]}, line width=1.3pt},
  drehung/.style={-{Stealth[length=1.8mm, width=1.6mm]}, line width=0.9pt},
  kraft/.style={-{Stealth[length=2.2mm, width=2mm]}, line width=1.1pt},
  mass/.style={{Stealth[length=1.4mm]}-{Stealth[length=1.4mm]}, line width=0.4pt, black!80},
  hilfslinie/.style={line width=0.3pt, black!60},
  marke/.style={circle, draw=black, fill=white, line width=0.8pt,
    minimum size=4.6mm, inner sep=0pt, font=\scriptsize\bfseries},
  zeiger/.style={line width=0.6pt},
  % Bewertung: \pic at (x,y) {ok}; bzw. {nok}
  ok/.pic={\fill[okgruen] (0,0) circle (3);
    \draw[white, line width=1.3pt, line cap=round, line join=round]
      (-1.4,0.1) -- (-0.4,-1.1) -- (1.5,1.2);},
  nok/.pic={\fill[nokrot] (0,0) circle (3);
    \draw[white, line width=1.3pt, line cap=round]
      (-1.1,-1.1) -- (1.1,1.1) (-1.1,1.1) -- (1.1,-1.1);},
}

% Feld mit hellem Hintergrund und Überschrift: \feld{x}{y}{Breite}{Höhe}{Titel}
\newcommand{\feld}[5]{%
  \fill[feld, rounded corners=2mm] (#1,#2) rectangle ++(#3,#4);
  \node[anchor=north west, font=\small\bfseries] at (#1+3,#2+#4-2) {#5};}

% Nummer an einer Zeichnung mit Zeigerlinie: \nr{Nummer}{Position}{Ziel}
\newcommand{\nr}[3]{%
  \draw[zeiger] (#2) -- (#3);
  \fill (#3) circle[radius=0.6];
  \node[marke] at (#2) {#1};}

% Kreissägeblatt in der Seitenansicht: \saegeblatt{x}{y}{Radius}{Zähne}
% Die Zähne zeigen im Uhrzeigersinn. In allen Seitenansichten steht die
% Bedienperson rechts und schiebt das Werkstück nach links: Die Zähne
% laufen oben auf sie zu und schneiden vorn von oben nach unten.
\newcommand{\saegeblatt}[4]{%
  \begin{scope}[shift={(#1,#2)}]
    \pgfmathsetmacro{\sbs}{360/#4}%
    \pgfmathsetmacro{\sbz}{min(0.18*#3, 2*pi*#3/#4*0.55)}%
    \fill[black!60, draw=black!85, line width=0.4pt] (0:#3)
      \foreach \i in {1,...,#4}{
        -- ({\i*\sbs-\sbs}:#3) -- ({\i*\sbs-0.2*\sbs}:#3-\sbz) -- ({\i*\sbs}:#3)}
      -- cycle;
    \fill[black!25] (0,0) circle (#3*0.5);
    \fill[white] (0,0) circle (#3*0.12);
  \end{scope}}
